\chapter{Recursive Definitions}
\signature

\section{The recursive idea}

A \emph{recursive definition} specifies an object in terms of smaller
instances of itself, anchored by explicit base cases. Three kinds of objects
are defined this way: \emph{functions} (a base value plus a recurrence),
\emph{sets} (base elements plus closure rules), and \emph{structures}
(atomic structures plus constructors). Induction (Chapter~13) is the proof
principle that matches each kind.

\section{Recursive functions}

\begin{definition}
A function on $\N$ is \emph{recursively defined} by giving $f(0)$ (and
possibly further initial values) and a rule computing $f(n+1)$ from earlier
values.
\end{definition}

\begin{example}
Factorial: $0! = 1$, $(n+1)! = (n+1) \cdot n!$.
Powers: $a^{0} = 1$, $a^{n+1} = a \cdot a^{n}$.
Fibonacci: $F_0 = 0$, $F_1 = 1$, $F_{n} = F_{n-1} + F_{n-2}$.
Ackermann's function --- recursive in two variables --- grows faster than
any primitive recursive function.
\end{example}

\subsection{From recurrence to closed form}

For \emph{linear homogeneous} recurrences, the characteristic-equation
method of Chapter~9 applies: $a_n = c_1 a_{n-1} + c_2 a_{n-2}$ has solution
$\alpha r_1^n + \beta r_2^n$ (distinct roots) or $(\alpha + \beta n)r^n$
(double root); Binet's formula
$F_n = \frac{\varphi^{n} - \psi^{n}}{\sqrt5}$ is the showcase. Simple
non-homogeneous patterns are handled by \emph{unrolling} (iteration) or
\emph{telescoping}.

\begin{example}[Telescoping]\label{ex:telescope}
$a_1 = 1$, $a_n = a_{n-1} + (2n - 1)$. Unroll:
$a_n = 1 + \sum_{k=2}^{n} (2k-1) = \sum_{k=1}^{n}(2k-1) = n^{2}$
(Chapter~13, Exercise 13.1).
\end{example}

\begin{example}[Tower of Hanoi]\label{ex:hanoi}
Moving $n$ discs requires $H_1 = 1$, $H_n = 2H_{n-1} + 1$ moves: move
$n - 1$ discs aside, the largest disc, then the $n-1$ again. Unrolling, or
induction, gives $H_n = 2^{n} - 1$. With $64$ discs:
$2^{64} - 1 \approx 1.8 \times 10^{19}$ moves --- the legend's monks will
not finish soon.
\end{example}

\section{Recursively defined sets}

\begin{definition}
A set $S$ is defined recursively by: a \emph{basis} (specified elements of
$S$), \emph{recursive rules} (ways to build new elements from old), and the
implicit \emph{exclusion} clause: nothing else is in $S$ ($S$ is the
\emph{smallest} set satisfying basis and rules).
\end{definition}

\begin{example}[Even naturals]
Basis: $0 \in E$. Rule: $x \in E \Rightarrow x + 2 \in E$. Then
$E = \{0, 2, 4, \ldots\}$, the even naturals.
\end{example}

\begin{example}[Strings]
The set $\Sigma^{*}$ of strings over alphabet $\Sigma$: basis --- the empty
string $\lambda \in \Sigma^{*}$; rule --- if $w \in \Sigma^{*}$ and
$a \in \Sigma$ then $wa \in \Sigma^{*}$. Concatenation, length and reversal
are then defined recursively on this structure, e.g.\
$\mathrm{len}(\lambda) = 0$, $\mathrm{len}(wa) = \mathrm{len}(w) + 1$.
\end{example}

\begin{example}
The well-formed formulas of Chapter~2 are a recursively defined set: atoms
as basis, connectives as rules --- which is what makes structural induction
on formulas legitimate.
\end{example}

\section{Recursively defined structures}

\begin{definition}
\emph{Rooted binary trees}: basis --- the empty tree; rule --- a root with
left and right binary subtrees forms a binary tree. \emph{Full} binary
trees: basis --- a single vertex; rule --- a root joined to two full binary
subtrees.
\end{definition}

\section{Structural induction}

\begin{theorem}[Structural induction]
To prove that every element of a recursively defined set has property $P$:
prove $P$ for each basis element, and prove that each rule preserves $P$
(if the inputs satisfy $P$, so does the output).
\end{theorem}

\begin{example}[Complete binary trees]\label{ex:tree}
The \emph{complete} full binary tree of height $n \geq 1$ has $2^{n} - 1$
vertices, of which $2^{n-1}$ are leaves. Induction on $n$: height $1$ ---
one vertex, one leaf: $2^1 - 1 = 1$, $2^{0} = 1$. Step: the tree of height
$n+1$ is a root over two trees of height $n$:
$2(2^{n} - 1) + 1 = 2^{n+1} - 1$ vertices and $2 \cdot 2^{n-1} = 2^{n}$
leaves.
\end{example}

\begin{example}[Strings]
$\mathrm{len}(uv) = \mathrm{len}(u) + \mathrm{len}(v)$: structural induction
on $v$. Base $v = \lambda$: $\mathrm{len}(u\lambda) = \mathrm{len}(u)$.
Rule case $v = wa$:
$\mathrm{len}(u(wa)) = \mathrm{len}((uw)a) = \mathrm{len}(uw) + 1
= \mathrm{len}(u) + \mathrm{len}(w) + 1 = \mathrm{len}(u) +
\mathrm{len}(wa)$.
\end{example}

\section{Recursion in computation}

\subsection{Recursive algorithms and correctness}

A recursive algorithm mirrors a recursive definition; induction proves it
correct. For the recursive factorial: base --- the program returns $1 = 0!$
on input $0$; step --- if it returns $n!$ on $n$, then on $n + 1$ it
returns $(n+1) \cdot n! = (n+1)!$. Mergesort splits a list in half, sorts
each recursively and merges; strong induction on length proves correctness,
and its running time obeys $T(n) = 2T(n/2) + O(n)$, solved by the
\emph{master theorem} as $T(n) = O(n \log n)$.

\subsection{Memoization and dynamic programming}

Naive recursive Fibonacci recomputes subproblems exponentially
($T(n) = T(n-1) + T(n-2) + O(1)$ is itself Fibonacci-sized). Caching each
computed value --- \emph{memoization} --- or filling a table bottom-up ---
\emph{dynamic programming} --- reduces the cost to linear. The lesson:
a recursive \emph{definition} need not dictate a naive recursive
\emph{computation}.

\subsection{Tail recursion, mutual recursion, fixed points}

A \emph{tail-recursive} function performs its recursive call last, so it
can be transformed mechanically into a loop (accumulator pattern):
$\mathrm{fact}(n, acc) = \mathrm{fact}(n-1, n \cdot acc)$. \emph{Mutual}
recursion defines functions in terms of each other, e.g.\
$\mathrm{even}(0) = \TT$, $\mathrm{even}(n) = \mathrm{odd}(n-1)$;
$\mathrm{odd}(0) = \FF$, $\mathrm{odd}(n) = \mathrm{even}(n-1)$ ---
proved correct by simultaneous induction. More abstractly, a recursive
definition presents its object as a \emph{fixed point} of an operator
($S = F(S)$); self-similar fractals such as the Koch curve and
Sierpi\'nski triangle, and the generations of a cellular automaton (e.g.\
rule 90, which paints Pascal's triangle mod 2), are recursive definitions
made visible.

\begin{notebox}
Common recursion pitfalls: a missing or unreachable base case (infinite
regress); recursive calls that do not strictly decrease toward the base;
overlapping subproblems solved exponentially often when a cache would do;
and confusing the definition of a set with membership testing --- the
exclusion clause matters.
\end{notebox}

\section{Summary}

Recursive definitions build functions (base values + recurrence), sets
(basis + closure rules + exclusion) and structures (atoms + constructors).
Closed forms come from characteristic equations, unrolling or telescoping;
structural induction is the universal proof method, with weak/strong
induction as its special case on $\N$. Computationally, recursion pairs
with induction for correctness, with memoization for efficiency, and with
iteration via tail calls.

\section*{Exercises}
\addcontentsline{toc}{section}{Exercises}

\begin{exercise}{\easy}
Give recursive definitions of: (a) $f(n) = 5n$; (b) $g(n) = 3^{n}$;
(c) the sequence $4, 9, 14, 19, \ldots$
\end{exercise}

\begin{exercise}{\easy}
Compute $a_4$ where $a_0 = 3$ and $a_{n} = 2a_{n-1} - 1$, then conjecture
and verify a closed form.
\end{exercise}

\begin{exercise}{\easy}
Define recursively the set $S = \{3, 6, 9, 12, \ldots\}$ of positive
multiples of $3$.
\end{exercise}

\begin{exercise}{\medium}
Solve by unrolling: $a_1 = 2$, $a_n = a_{n-1} + 2^{n}$, and verify by
induction.
\end{exercise}

\begin{exercise}{\medium}
Prove by induction that the Tower of Hanoi recurrence $H_1 = 1$,
$H_n = 2H_{n-1} + 1$ has solution $H_n = 2^{n} - 1$.
\end{exercise}

\begin{exercise}{\medium}
Define the reversal of a string by $\lambda^{R} = \lambda$ and
$(wa)^{R} = a\,w^{R}$. Prove by structural induction that
$(uv)^{R} = v^{R} u^{R}$.
\end{exercise}

\begin{exercise}{\medium}
Let $S$ be defined by: $5 \in S$; if $x, y \in S$ then $x + y \in S$.
Determine $S$ exactly and prove your answer by structural induction (one
inclusion) and ordinary argument (the other).
\end{exercise}

\begin{exercise}{\medium}
A \emph{full} binary tree is a single vertex, or a root with two full
binary subtrees. Prove by structural induction that the number of leaves
of a full binary tree exceeds the number of internal vertices by exactly
one.
\end{exercise}

\begin{exercise}{\hard}
Define $T(n) = T(\lfloor n/2 \rfloor) + 1$, $T(1) = 1$. Show
$T(n) = \lfloor \log_2 n \rfloor + 1$ by strong induction.
\end{exercise}

\begin{exercise}{\hard}
Naive recursive Fibonacci makes $C_n$ function calls, where $C_0 = C_1 = 1$
and $C_n = C_{n-1} + C_{n-2} + 1$. Prove $C_n = 2F_{n+1} - 1$, confirming
exponential blow-up.
\end{exercise}

\begin{exercise}{\hard}
The Koch curve is defined recursively: $K_0$ is a segment of length $1$;
$K_{n+1}$ replaces each segment of $K_n$ by four segments each one third as
long. Find the length $L_n$ of $K_n$, prove it by induction, and conclude
about $\lim L_n$.
\end{exercise}

\section*{Solutions to the Exercises}
\addcontentsline{toc}{section}{Solutions to the Exercises}

\begin{solution}{14.1}
(a) $f(0) = 0$, $f(n+1) = f(n) + 5$.
(b) $g(0) = 1$, $g(n+1) = 3\,g(n)$.
(c) $a_0 = 4$, $a_{n+1} = a_n + 5$.
\end{solution}

\begin{solution}{14.2}
$a_1 = 5$, $a_2 = 9$, $a_3 = 17$, $a_4 = 33$. Conjecture
$a_n = 2^{n+1} + 1$: base $a_0 = 3$; step
$2(2^{n+1} + 1) - 1 = 2^{n+2} + 1$. \checkmark
\end{solution}

\begin{solution}{14.3}
Basis: $3 \in S$. Rule: $x \in S \Rightarrow x + 3 \in S$. (Exclusion
clause: nothing else.) An easy induction shows $S = \{3k \mid k \ge 1\}$.
\end{solution}

\begin{solution}{14.4}
$a_n = 2 + \sum_{k=2}^{n} 2^{k} = 2 + (2^{n+1} - 4) = 2^{n+1} - 2$.
Induction: base $a_1 = 2 = 2^2 - 2$; step
$a_{n+1} = (2^{n+1} - 2) + 2^{n+1} = 2^{n+2} - 2$.
\end{solution}

\begin{solution}{14.5}
Base: $H_1 = 1 = 2^1 - 1$. Step:
$H_{n+1} = 2H_n + 1 = 2(2^{n} - 1) + 1 = 2^{n+1} - 1$.
\end{solution}

\begin{solution}{14.6}
Structural induction on $v$. Base $v = \lambda$:
$(u\lambda)^R = u^R = \lambda^R u^R$. Step $v = wa$:
$(u(wa))^R = ((uw)a)^R = a\,(uw)^R = a\,(w^R u^R)
= (a\,w^R)\,u^R = (wa)^R u^R$.
\end{solution}

\begin{solution}{14.7}
$S = \{5k \mid k \geq 1\}$, the positive multiples of $5$. ($\subseteq$)
Structural induction: the basis element $5$ is a multiple of $5$, and the
sum of two multiples of $5$ is one. ($\supseteq$) Induction on $k$:
$5 \in S$; if $5k \in S$ then $5(k+1) = 5k + 5 \in S$ by the rule.
\end{solution}

\begin{solution}{14.8}
$P(T)$: leaves$(T)$ = internal$(T) + 1$. Base: a single vertex has $1$ leaf,
$0$ internal: \checkmark. Step: $T$ = root over $T_1, T_2$ satisfying $P$.
Then leaves$(T)$ = leaves$(T_1)$ + leaves$(T_2)$ =
internal$(T_1)$ + internal$(T_2)$ + 2 = internal$(T)$ + 1, since
internal$(T)$ = internal$(T_1)$ + internal$(T_2)$ + 1 (the new root).
\end{solution}

\begin{solution}{14.9}
$P(n)$: $T(n) = \lfloor \log_2 n \rfloor + 1$. Base: $T(1) = 1 = 0 + 1$.
Step: let $n \geq 2$, so $1 \le \lfloor n/2 \rfloor < n$; by strong
hypothesis $T(n) = \lfloor \log_2 \lfloor n/2 \rfloor \rfloor + 2$. Writing
$2^{k} \le n < 2^{k+1}$ (so $\lfloor\log_2 n\rfloor = k$), we get
$2^{k-1} \le \lfloor n/2 \rfloor < 2^{k}$, whence
$\lfloor \log_2 \lfloor n/2 \rfloor \rfloor = k - 1$ and
$T(n) = (k-1) + 2 = k + 1$. \checkmark
\end{solution}

\begin{solution}{14.10}
Induction with two bases: $C_0 = 1 = 2F_1 - 1$, $C_1 = 1 = 2F_2 - 1$.
Step:
$C_{n} = C_{n-1} + C_{n-2} + 1
= (2F_{n} - 1) + (2F_{n-1} - 1) + 1 = 2F_{n+1} - 1$.
Since $F_{n+1} \sim \varphi^{n+1}/\sqrt5$, the call count grows
exponentially in $n$.
\end{solution}

\begin{solution}{14.11}
Each step multiplies the number of segments by $4$ and divides their length
by $3$, so $L_{n} = (4/3)^{n}$. Induction: $L_0 = 1$; step
$L_{n+1} = 4 \cdot \frac{L_n}{3} \cdot \frac{\text{(per
segment)}}{\text{(same)}} = \tfrac43 L_n = (4/3)^{n+1}$. As
$4/3 > 1$, $L_n \to \infty$: the limiting curve has infinite length
(while spanning a bounded region --- the hallmark of a fractal).
\end{solution}
